% =============================================================================
%  3.2 Composition API
% =============================================================================
\section{Composition API}

\term{Composition API}{Composition API} — набор функций, позволяющий
организовывать логику компонента по функциональным блокам, а не по
опциям (\texttt{data}, \texttt{methods}, \texttt{computed}).  Это решает
проблему «разбросанной логики» в крупных компонентах, написанных
с помощью Options API.

Composition API стал основным рекомендуемым подходом начиная с Vue~3.
Он позволяет извлекать повторно используемую логику в
\term{композабл-функции}{Composables}.

\subsection{setup и script setup}

\begin{code}{html}
<script setup>
import { ref, computed } from 'vue';

const items = ref([]);
const count = computed(() => items.value.length);

function addItem(item) {
  items.value.push(item);
}
</script>

<template>
  <p>Элементов: {{ count }}</p>
  <button @click="addItem('новый')">Добавить</button>
</template>
\end{code}

\subsection{Composables}

\subsubsection{Паттерн извлечения логики}

Композабл — это функция, начинающаяся с \texttt{use}, которая
инкапсулирует реактивное состояние и логику.  Количество повторного
использования $R$ композабла определяется числом компонентов, которые
его импортируют: при $R = n$ экономия строк составляет приблизительно
$(n - 1) \cdot L$, где $L$ — число строк логики.

\begin{equation}
  \text{Saved LOC} \approx (n - 1) \cdot L
\end{equation}

\begin{code}{javascript}
// composables/useFetch.js
import { ref, watchEffect } from 'vue';

export function useFetch(url) {
  const data = ref(null);
  const error = ref(null);

  watchEffect(async () => {
    try {
      const res = await fetch(url.value);
      data.value = await res.json();
    } catch (e) {
      error.value = e;
    }
  });

  return { data, error };
}
\end{code}

\note{Композабл-функции должны вызываться синхронно внутри
\texttt{setup()} или \texttt{<script setup>}.  Это позволяет Vue
корректно привязать реактивные зависимости к текущему экземпляру
компонента.}

\important{Избегайте создания побочных эффектов в
\texttt{computed()} — вычисляемые свойства должны быть чистыми
функциями.  Для побочных эффектов используйте \texttt{watch()}
или \texttt{watchEffect()}.}
